\SixthArticle{24}
Since the limit of the function $F(\alpha,\beta,\gamma+k,1)$, as $k$ grows to infinity, is evidently unity, equation 39 becomes
\[
\begin{array}{@{}ll@{}}
[48]&\displaystyle F(\alpha,\beta,\gamma,1)=
\frac{\Pi(\gamma-1).\Pi(\gamma-\alpha-\beta-1)}
{\Pi(\gamma-\alpha-1).\Pi(\gamma-\beta-1)}
\end{array}
\]
This formula gives the complete solution of the problem that forms the subject of this section. The elegant equations below follow at once:
\SourceLeaf{157}
\[
\begin{array}{@{}ll@{}}
[49]&\displaystyle F(\alpha,\beta,\gamma,1)=F(-\alpha,-\beta,\gamma-\alpha-\beta,1)\\[3pt]
[50]&\displaystyle F(\alpha,\beta,\gamma,1).F(-\alpha,\beta,\gamma-\alpha,1)=1\\[3pt]
[51]&\displaystyle F(\alpha,\beta,\gamma,1).F(\alpha,-\beta,\gamma-\beta,1)=1
\end{array}
\]
where $\gamma$ in the first, $\gamma-\beta$ in the second, and $\gamma-\alpha$ in the third must be a positive quantity.

\SixthArticle{25}
Let us apply formula 48 to some of the equations of Article 5. On setting $t=90^0=\frac12\pi$, formula XIII becomes $\frac12\pi=F(\frac12,\frac12,\frac32,1)$, or is equivalent to the known equation
\[
\frac12\pi=1+\frac{1.1}{2.3}+\frac{1.1.3}{2.4.5}
+\frac{1.1.3.5}{2.4.6.7}+\mathrm{etc.}
\]
Thus, since formula 48 gives $F(\frac12,\frac12,\frac32,1)=\frac{\Pi\frac12.\Pi(-\frac12)}{\Pi0.\Pi0}$, and since $\Pi0=1$, $\Pi\frac12=\frac12\Pi(-\frac12)$, it follows that $\pi=(\Pi(-\frac12))^2$, or
\[
\begin{array}{@{}ll@{}}
[52]&\displaystyle\Pi(-\tfrac12)=\BareRoot{\pi}\\[3pt]
[53]&\displaystyle\Pi\tfrac12=\tfrac12\BareRoot{\pi}
\end{array}
\]
Formula XVI of Article 5, which is equivalent to the known equation
\[
\sin nt=n\sin t-\frac{n(nn-1)}{2.3}\sin t^3
+\frac{n(nn-1)(nn-9)}{2.3.4.5}\sin t^5-\mathrm{etc.}
\]
and holds in general for any value of $n$, provided $t$ remains within the limits $-90^0$ and $+90^0$, gives for $t=\frac12\pi$
\[
\sin\frac{n\pi}{2}=
\frac{n\Pi\frac12.\Pi(-\frac12)}{\Pi(-\frac12n).\Pi\frac12n}
\]
from which the elegant formula follows:
\[
\Pi\frac12n.\Pi(-\frac12n)=\frac{\frac12n\pi}{\sin\frac12n\pi},
\quad\hbox{or, on setting }n=2z
\]
\[
\begin{array}{@{}ll@{}}
[54]&\displaystyle\Pi(-z).\Pi(+z)=\frac{z\pi}{\sin z\pi}\\[3pt]
[55]&\displaystyle\Pi(-z).\Pi(z-1)=\frac{\pi}{\sin z\pi}
\end{array}
\]
and also, on writing $z+\frac12$ for $z$,
\[
\begin{array}{@{}ll@{}}
[56]&\displaystyle\Pi(-\tfrac12+z).\Pi(-\tfrac12-z)=\frac{\pi}{\cos z\pi}
\end{array}
\]
\SourceLeaf{158}
Combining formula 54 with the definition of the function $\Pi$ shows that $\frac{z\pi}{\sin z\pi}$ is the limit of the infinite product
\[
\frac{(1.2.3.4\,.\,.\,.\,.k)^2}
{(1-zz)(4-zz)(9-zz)\,.\,.\,.\,.(kk-zz)}
\]
as $k$ grows to infinity, and hence
\[
\sin z\pi=z\pi(1-zz)(1-\frac{zz}{4})(1-\frac{zz}{9})\quad
\mathrm{etc.\ to\ infinity}
\]
Similarly, from 56 follows
\[
\cos z\pi=(1-4zz)(1-\frac{4zz}{9})(1-\frac{4zz}{25})\quad
\mathrm{etc.\ to\ infinity}
\]
These are very well-known formulas, which analysts usually derive by entirely different methods.

\SixthArticle{26}
With $n$ denoting an integer, the value of the expression
\[
\frac{n^{nz}\Pi(k,z).\Pi(k,z-\frac1n).\Pi(k,z-\frac2n)
\,.\,.\,.\,.\Pi(k,z-\frac{n-1}{n})}{\Pi(nk,nz)}
\]
when properly collected is found to be
\[
=\frac{(1.2.3\,.\,.\,.\,.k)^n n^{nk}}
{1.2.3\,.\,.\,.\,.nk.k^{\frac12(n-1)}}
\]
and is therefore independent of $z$, or will remain the same whatever value is assigned to $z$. Thus, since $\Pi(k,0)=\Pi(nk,0)=1$, it can be represented by the product
\[
\Pi(k,-\frac1n).\Pi(k,-\frac2n).\Pi(k,-\frac3n)
\,.\,.\,.\,.\Pi(k,-\frac{n-1}{n})
\]
Hence, as $k$ grows to infinity, we obtain
\[
\begin{gathered}
\frac{n^{nz}\Pi z.\Pi(z-\frac1n).\Pi(z-\frac2n)
\,.\,.\,.\,.\Pi(z-\frac{n-1}{n})}{\Pi nz}\\[4pt]
=\Pi(-\frac1n).\Pi(-\frac2n).\Pi(-\frac3n)
\,.\,.\,.\,.\Pi(-\frac{n-1}{n})
\end{gathered}
\]
The product on the right, multiplied by itself with the order of its factors reversed, gives, by formula 55,
\[
\frac{\pi}{\sin\frac1n\pi}.\frac{\pi}{\sin\frac2n\pi}.
\frac{\pi}{\sin\frac3n\pi}\,.\,.\,.\,.
\frac{\pi}{\sin\frac{n-1}{n}\pi}=\frac{(2\pi)^{n-1}}{n}
\]
Thus we have the elegant theorem
\SourceLeaf{159}
\[
\begin{aligned}
{}\text{[57]}\quad &
\frac{n^{nz}\Pi z.\Pi(z-\frac1n).\Pi(z-\frac2n).\,.\,.\,.
\Pi(z-\frac{n-1}{n})}{\Pi nz}
=\frac{(2\pi)^{\frac12(n-1)}}{\BareRoot n}
\end{aligned}
\]

